
\section{\trReducedness{}与\trIntegrality{}}

回想一下，概形令人警觉的事情之一便是：函数并不由其在各点处的值决定，而这正是%
\trNilpotent{}元在作祟（\S3.2.12）。

\textbf{5.2.1.}\ \emph{\trhDefinition{}。}\quad 我们称一个\trScheme{} $X$ 是%
\textbf{\trReduced}（reduced）的，如果它的诸\trStalk{} $\OO_{X,p}$ 都没有非零%
\trNilpotent{}元（对一切 $p$）。\trReducedness{}按定义就是一条\trStalkLocal{}性质（\S2.4.1）。

\textbf{5.2.A.（练习）}\quad 证明：若 $f$ 与 $g$ 是\trReduced{}\trScheme{}上的两个函数（即 $\OO_X$ 的两个%
\trGlobalSections），二者在所有点处取值都相同，则 $f=g$。（两条提示：单射%
$\OO_X(\mathsf{U})\hookrightarrow\prod_{p\in\mathsf{U}}\OO_{X,p}$，来自%
\cref{ex:sections-determined-by-germs}；以及\trNilradical{}是全体\trPrimeIdeal{}的交，
来自\cref{thm:nilradical-intersection-primes}。）

\textbf{5.2.B.（练习）}\quad 回想：若一个\trRing{}没有非零\trNilpotent{}元，就称它是\emph{\trReduced}的（\S3.2.14）。
证明：$A$ 是\trReduced{}环，当且仅当 $\Spec A$ 是\trReduced{}的。再证明：若对 $X$ 的
每个\trOpenSet{} $\mathsf{U}$，$\OO_X(\mathsf{U})$ 都是\trReduced{}环，则\trScheme{}
$X$ 是\trReduced{}的。

\textbf{5.2.C.（容易的练习）}\quad 证明 $\mathbb{A}^n_k$ 与 $\mathbb{P}^n_k$ 都是\trReduced{}的。

\trScheme{} $\Spec k[x,y]/(y^2,xy)$ 是非\trReduced{}的。在图~\ref{fig:nonreduced-y2xy}
中勾画它时，我们曾示意：那团\trFuzz{}标示出的正是原点处的非\trReducedness{}。下面的
练习便是把这件事说精确的第一次尝试。

\textbf{5.2.D.（练习）}\quad 证明 $\bigl(k[x,y]/(y^2,xy)\bigr)_x$ 没有非零\trNilpotent{}元。（可能的提示：借助几何
图画，证明它\trIsomorphic{}于另一个\trRing。\cref{ex:xy-localization-isomorphism} 也许
能给出另一条提示。）再证明：$\Spec k[x,y]/(y^2,xy)$ 中唯一具有非\trReduced{}\trStalk{}
的点就是原点。

\textbf{5.2.E.（不重要的练习）}\quad 设 $X$ 是\trQuasicompact{}\trScheme，证明：只需在\trClosedPoint{}处检验\trReducedness{}。提示：
别试图证明\trReducedness{}是一条\textbf{开条件}（open condition，见注记~5.2.2）。作为替代，利用
练习~5.1.E 证明：任何非\trReduced{}的点，其闭包中都有一个非\trReduced{}的%
\trClosedPoint。（这个结果很有意思，但我们用不上它。）

\textbf{5.2.2.}\ \emph{\trhRemark{}。}\quad 我们将在练习~6.6.N 中看到：对%
\trLocally\trNoetherian{}\trScheme，\trReducedness{}是一条开条件。但一般地，\trReducedness{}并不是开条件。
你或许能把
\[
X=\Spec\mathbb{C}[x,y_1,y_2,\dots]\big/\bigl(y_1^2,\,y_2^2,\,y_3^2,\,\dots,\,(x-1)y_1,\,(x-2)y_2,\,(x-3)y_3,\,\cdots\bigr)
\]
的底\trTopologicalSpace{}与 $\Spec\mathbb{C}[x]$ 的底\trTopologicalSpace{}认同起来，
进而证明：$X$ 的非\trReduced{}点恰是与正整数相对应的那些\trClosedPoint。此集合的%
\trComplement{}不是\trZariski{}开集。

\textbf{5.2.3.}\ \hard\ \emph{给专家的\trhWarning{}。}\quad 若\trScheme{} $X$ 是%
\trReduced{}的，则（据练习~5.2.B）其\trGlobalSections{}所成之%
\trRing{}是\trReduced{}的。然而逆命题不真。你现在已经有足够的知识来验证：若 $X$ 是
$\mathbb{P}^2_k$ 中由 $x^2=0$ 割出的\trScheme，则 $\Gamma(X,\OO_X)\cong k$，而 $X$ 非%
\trReduced{}。（这个例子日后还将在 \S11.4.7 与 \S18.6.U 中出现。）

\textbf{5.2.F.（练习）}\quad 设 $X$ 是\trQuasicompact{}的，而 $f$ 是一个在 $X$ 的所有点处都消失的函数。证明：存在
某个 $n$ 使得 $f^n=0$。再证明：当 $X$ 不\trQuasicompact{}时，这可能失效。（本练习较不
重要，但它说明了我们为何钟爱\trQuasicompactness，并给出了\trQuasicompactness{}不成立
时的一种标准病态。）提示：取诸 $\Spec A_n$ 的一个无穷\trDisjointUnion，其中%
$A_n:=k[\epsilon]/(\epsilon^n)$。这个\trScheme{}日后还将在 \S9.4.2（其图画见图~9.5）
以及警示/例~9.4.11 中再次出现。

\textbf{5.2.4.}\ \emph{\trhDefinition{}。}\quad \trScheme{} $X$ 称为%
\textbf{\trIntegral}（integral）的，如果它非空，并且对 $X$ 的每个非空\trOpenSubset{}%
$\mathsf{U}$，$\OO_X(\mathsf{U})$ 都是\trIntegralDomain。

\textbf{5.2.G.（重要练习）}\quad 证明：\trScheme{} $X$ 是\trIntegral{}的，当且仅当它既\trIrreducible{}又\trReduced{}。
（于是我们把\trIntegral{}\trScheme{}画成：「一整块，不惹\trFuzz{}。」可能的提示：%
\cref{ex:vanishing-locus-closed}。）

\textbf{5.2.H.（练习）}\quad 证明：\trAffineScheme{} $\Spec A$ 是\trIntegral{}的，当且仅当 $A$ 是\trIntegralDomain。

\textbf{5.2.I.（练习）}\quad 设 $X$ 是\trIntegral{}\trScheme。则 $X$（因为它\trIrreducible）有一个\trGenericPoint{}%
$\eta$。设 $\Spec A$ 是 $X$ 的任一非空\trAffine\trOpenSubset。证明：$\OO_{X,\eta}$%
（即 $\OO_X$ 在 $\eta$ 处的\trStalk）被自然地认同于 $A$ 的\trFractionField{}%
$\mathrm{K}(A)$。它称为 $X$ 的\textbf{\trFunctionField}（function field）$\mathrm{K}(X)$。
在 $X$ 的任何非空\trOpenSet{}上都能把它算出来，因为任何这样的\trOpenSet{}都含有%
\trGenericPoint。$\mathrm{K}(X)$ 的元素称为 $X$ 上的\textbf{\trRationalFunction}
（rational function；日后将在定义~6.6.36 中进一步推广）。

\textbf{5.2.J.（练习）}\quad 设 $X$ 是\trIntegral{}\trScheme。证明：只要 $\mathsf{V}\neq\varnothing$，诸%
\trRestrictionMap{} $\operatorname{res}_{\mathsf{U},\mathsf{V}}\colon\OO_X(\mathsf{U})
\to\OO_X(\mathsf{V})$ 就都是包含。设 $\Spec A$ 是 $X$ 的任一非空\trAffine\trOpenSubset%
（于是 $A$ 是\trIntegralDomain）。证明：自然映射
$\OO_X(\mathsf{U})\to\OO_{X,\eta}=\mathrm{K}(A)$（其中 $\mathsf{U}$ 是任一非空%
\trOpenSubset，而 $\eta$ 是\trGenericPoint）是包含。

于是，不可约\trVariety（日后定义的\trIntegral{}\trScheme{}的一个重要例子）便享有一件
方便的性质：不同\trOpenSet{}上的\trSection{}可以被看作同一个\trRing{}的子集。特别地，
\trRestrictionMap（除指向空集的情形外）总是包含，而粘合也变得容易：$\mathsf{U}$ 的一个%
\trCover{} $\mathsf{U}_i$ 上的诸函数 $f_i$（$i$ 跑遍一个\trIndexSet）能够粘合，当且仅当
它们是 $\mathrm{K}(X)$ 的同一个元素。这是（不可约的）\trVariety{}通常先于\trScheme{}被
引入的原因之一。

\textbf{5.2.5.}\ \emph{\trhCaveat{}。}\quad \trIntegrality{}不是\trStalkLocal{}性质（两个%
\trIntegral{}\trScheme{}的\trDisjointUnion{}不是\trIntegral{}的，因为据%
\cref{ex:product-ring-homeomorphism}，$\Spec A\sqcup\Spec B=\Spec(A\times B)$），但它
差一点就是；见练习~5.3.C。
